\documentclass[11pt,a4paper]{article}

\usepackage[T1]{fontenc}
\usepackage[utf8]{inputenc}
\usepackage[english]{babel}
\usepackage{lmodern}
\usepackage{microtype}
\usepackage{amsmath,amssymb,amsthm,mathtools}
\usepackage{array}
\usepackage{longtable}
\usepackage[a4paper,left=25mm,right=25mm,top=25mm,bottom=25mm]{geometry}
\usepackage[hidelinks]{hyperref}

\hypersetup{
  pdftitle={The Components of the Powers of z = x + iy + jt: Coefficients, Recursion and the Laplace-Beltrami Equation},
  pdfkeywords={Reduced quaternions, Powers, Real part, Binomial coefficients, Laplace-Beltrami, Lean}
}

\newcommand{\ii}{\mathrm{i}}
\newcommand{\jj}{\mathrm{j}}
\DeclareMathOperator{\Real}{Re}
\newcommand{\R}{\mathbb{R}}
\newcommand{\N}{\mathbb{N}}
\newcommand{\HH}{\mathbb{H}}
\newcommand{\lean}[1]{\texttt{#1}}
\DeclarePairedDelimiter{\abs}{\lvert}{\rvert}

\newtheorem{theorem}{Theorem}
\newtheorem{corollary}[theorem]{Corollary}
\newtheorem{lemma}[theorem]{Lemma}
\theoremstyle{remark}
\newtheorem{remark}[theorem]{Remark}
\newtheorem{example}[theorem]{Example}

\title{The Components of the Powers of $z=x+\ii y+\jj t$:\\
Coefficients, Recursion and the Laplace--Beltrami Equation}
\author{}
\date{}

\begin{document}
\maketitle

\begin{abstract}
For the reduced quaternion $z=x+\ii y+\jj t$ we determine $\Real(z^n)$ explicitly,
give the closed form and the two-term recursion of its coefficients, and settle the
conjectures of the accompanying handwritten notes:
the coefficient symmetry holds for every index, and $t\,\Delta u=\partial u/\partial t$
holds for every $n$; in fact this equation is \emph{equivalent} to the recursion.
The same is done for the coefficients of $\ii$ and $\jj$: $z^n=\Real(z^n)+(\ii y+\jj t)S_n$
with explicit $S_n$, the coefficient of $\ii$ satisfies the same equation, and the coefficient
of $\jj$ satisfies $t^2\Delta w-t\,\partial_tw+w=0$.
Finally everything is extended to paravectors $x+\sum x_ie_i$ of the Clifford algebra
$\mathrm{Cl}_{0,\mu}$: the polynomials are the same and the equations become
$x_l\Delta f=(\mu-1)\partial_{x_l}f$ and
$x_l^2\Delta w-(\mu-1)x_l\partial_{x_l}w+(\mu-1)w=0$.
Every statement is machine-checked in Lean~4 with Mathlib.
\end{abstract}

\section{Setting and results}

Let $\HH$ be the real quaternions with $\ii^2=\jj^2=-1$, $\ii\jj=-\jj\ii$, and
\[
  z=x+\ii y+\jj t,\qquad v=\ii y+\jj t,\qquad r=y^2+t^2,\qquad (x,y,t)\in\R^3 .
\]

\begin{theorem}[Closed form]\label{thm:closed}
For every $n\in\N$,
\[
  \Real(z^n)=\sum_{k=0}^{\lfloor n/2\rfloor} c_k\,x^{n-2k}\,r^k,
  \qquad c_k=c_k^{(n)}=(-1)^k\binom{n}{2k}.
\]
In particular, for $n=2m$ the coefficients of the notes are
\[
  \boxed{\;a_k=(-1)^k\binom{2m}{2k},\qquad 0\le k\le m.\;}
\]
\end{theorem}

\begin{proof}
We have $v^2=-y^2-t^2+(\ii\jj+\jj\ii)yt=-r\in\R$, and $x\in\R$ is central, so the
binomial theorem applies to $z=x+v$:
\[
  z^n=\sum_{i=0}^{n}\binom{n}{i}x^{n-i}v^{i}.
\]
Now $v^{2k}=(-r)^k\in\R$ and $v^{2k+1}=(-r)^k v$ with $\Real v=0$. Taking real parts
kills the odd terms and leaves the asserted sum.
\end{proof}

No polar representation is used; the axis $y=t=0$ requires no separate treatment.

\begin{corollary}[Recursion]\label{cor:rec}
With $c_0=1$, for all $k\ge 0$,
\[
  (2k+2)(2k+1)\,c_{k+1}=-(n-2k)(n-2k-1)\,c_k ,
\]
that is, for $n=2m$,
\[
  \boxed{\;a_0=1,\qquad
  a_{k+1}=-\,\frac{(2m-2k)(2m-2k-1)}{(2k+2)(2k+1)}\,a_k .\;}
\]
\end{corollary}

\begin{proof}
Apply $(j+1)\binom{n}{j+1}=(n-j)\binom{n}{j}$ with $j=2k+1$ and $j=2k$. For
$2k\ge n$ both sides vanish.
\end{proof}

\begin{corollary}[Symmetry]\label{cor:sym}
For $n=2m$ and every $0\le k\le m$,
\[
  a_{m-k}=(-1)^m a_k ,\qquad\text{in particular } a_m=(-1)^m .
\]
\end{corollary}

\begin{proof}
$\binom{2m}{2m-2k}=\binom{2m}{2k}$ and $(-1)^{m-k}=(-1)^m(-1)^k$.
\end{proof}

\begin{remark}
The notes conjecture $a_{m-(2k-1)}=(-1)^m a_{2k-1}$ for $m\ge 2k+1$.
Corollary~\ref{cor:sym} shows that neither the parity of the index nor the lower
bound on $m$ is needed.
\end{remark}

\section{The Laplace--Beltrami equation}

Let $\Delta=\partial_x^2+\partial_y^2+\partial_t^2$. The equation
$t\,\Delta u=\partial_t u$, i.e.\ $t\,\Delta u-\partial_t u=0$, is the
Laplace--Beltrami equation of the hyperbolic metric $t^{-2}(dx^2+dy^2+dt^2)$ on
the upper half space, up to the factor $t$.

\begin{lemma}\label{lem:axial}
Let $P\in\R[x,r]$ and $u(x,y,t)=P(x,y^2+t^2)$. Then
\[
  t\,\Delta u-\partial_t u=t\,\bigl(P_{xx}+4rP_{rr}+2P_r\bigr)\big|_{r=y^2+t^2}.
\]
\end{lemma}

\begin{proof}
$\partial_t u=2tP_r$, $\partial_t^2u=2P_r+4t^2P_{rr}$, likewise in $y$; hence
$\Delta u=P_{xx}+4P_r+4rP_{rr}$ and $t\,\Delta u-\partial_tu=t(P_{xx}+4rP_{rr}+2P_r)$.
\end{proof}

\begin{theorem}\label{thm:lb}
Let $n\in\N$, let $c_0,\dots,c_{\lfloor n/2\rfloor}\in\R$, put $c_k=0$ for
$k>\lfloor n/2\rfloor$, and
\[
  u=\sum_{k}c_k\,x^{n-2k}(y^2+t^2)^k .
\]
Then
\[
  t\,\Delta u=\frac{\partial u}{\partial t}
  \quad\Longleftrightarrow\quad
  (2k+2)(2k+1)\,c_{k+1}=-(n-2k)(n-2k-1)\,c_k\ \text{ for all } k\ge 0 .
\]
Consequently $u_n=\Real(z^n)$ satisfies $t\,\Delta u_n=\partial_t u_n$ for every
$n\in\N$, and $u_n$ is the unique solution of this form with $c_0=1$.
\end{theorem}

\begin{proof}
With $P=\sum_k c_kx^{n-2k}r^k$ one computes
\[
  P_{xx}+4rP_{rr}+2P_r
  =\sum_{k\ge 0}\Bigl[(n-2k)(n-2k-1)\,c_k+(2k+2)(2k+1)\,c_{k+1}\Bigr]x^{n-2k-2}r^k,
\]
using $4k(k-1)+2k=2k(2k-1)$ and an index shift. Since the map
$(x,y,t)\mapsto(x,y^2+t^2)$ has image with non-empty interior, a polynomial in
$(x,r)$ vanishing on it is zero; by Lemma~\ref{lem:axial} the equation therefore
holds iff every bracket vanishes. The second assertion follows from
Corollary~\ref{cor:rec}.
\end{proof}

\begin{remark}
Conceptually: $P(x,\rho^2)=\Real(x+\ii\rho)^n$ is harmonic in the $(x,\rho)$-plane,
and $P_{xx}+4rP_{rr}+2P_r=P_{xx}+P_{\rho\rho}$ for $r=\rho^2$. Thus
Theorem~\ref{thm:lb} says that the axially symmetric lift of a plane harmonic
function $f(x,\rho)$, even in $\rho$, satisfies $\Delta u=\rho^{-1}f_\rho$ and
$\partial_tu=t\rho^{-1}f_\rho$. The same argument applies to $\Real$ of any power
series in $z$ with real coefficients.
\end{remark}

\section{The components of $\ii$ and $\jj$}

The expansion used for Theorem~\ref{thm:closed} also yields the remaining components.

\begin{theorem}[Components of $\ii$ and $\jj$]\label{thm:imag}
For every $n\in\N$,
\[
  \boxed{\;z^n=\Real(z^n)+(\ii y+\jj t)\,S_n,\qquad
  S_n=\sum_{k=0}^{\lfloor (n-1)/2\rfloor}s_k\,x^{n-1-2k}\,r^k,\qquad
  s_k=s_k^{(n)}=(-1)^k\binom{n}{2k+1}.\;}
\]
Thus the coefficient of $\ii$ is $v_n=y\,S_n$, the coefficient of $\jj$ is $w_n=t\,S_n$, and
the coefficient of $\ii\jj$ vanishes: the powers of a reduced quaternion are reduced.
\end{theorem}

\begin{proof}
In $z^n=\sum_i\binom{n}{i}x^{n-i}v^i$ the even powers $v^{2k}=(-r)^k$ are real and the odd
powers are $v^{2k+1}=(-r)^k v$ with $v=\ii y+\jj t$.
\end{proof}

For instance $z^2=x^2-r+2x\,v$, $z^3=x^3-3xr+(3x^2-r)\,v$ and
$z^4=x^4-6x^2r+r^2+(4x^3-4xr)\,v$.

\begin{corollary}[Recursion and symmetry]\label{cor:srec}
With $s_0=n$, for all $k\ge 0$,
\[
  \boxed{\;s_{k+1}=-\,\frac{(n-2k-1)(n-2k-2)}{(2k+3)(2k+2)}\,s_k ,\;}
\]
and for $n=2m$ and every $0\le k\le m-1$,
\[
  s_{m-1-k}=(-1)^{m-1}s_k .
\]
\end{corollary}

\begin{proof}
Apply $(j+1)\binom{n}{j+1}=(n-j)\binom{n}{j}$ with $j=2k+2$ and $j=2k+1$; for $2k+3>n$
both sides vanish. The symmetry is $\binom{2m}{2m-2k-1}=\binom{2m}{2k+1}$.
\end{proof}

\begin{lemma}\label{lem:axialim}
Let $P\in\R[x,r]$, $v(x,y,t)=y\,P(x,y^2+t^2)$ and $w(x,y,t)=t\,P(x,y^2+t^2)$. Then
\begin{align*}
  t\,\Delta v-\partial_t v&=t\,y\,\bigl(P_{xx}+4rP_{rr}+6P_r\bigr)\big|_{r=y^2+t^2},\\
  t^2\Delta w-t\,\partial_t w+w&=t^3\bigl(P_{xx}+4rP_{rr}+6P_r\bigr)\big|_{r=y^2+t^2}.
\end{align*}
\end{lemma}

\begin{proof}
With $u=P(x,y^2+t^2)$ one has $\Delta(yu)=y\,\Delta u+2\,\partial_yu$ and
$\Delta(tu)=t\,\Delta u+2\,\partial_tu$, where $\partial_yu=2yP_r$, $\partial_tu=2tP_r$ and,
by Lemma~\ref{lem:axial}, $\Delta u=P_{xx}+4rP_{rr}+4P_r$. Hence
$t\,\Delta v-\partial_tv=ty(P_{xx}+4rP_{rr}+4P_r)+4tyP_r-2tyP_r$ and
$t^2\Delta w-t\,\partial_tw+w=t^3(P_{xx}+4rP_{rr}+4P_r)+4t^3P_r-tu-2t^3P_r+tu$.
\end{proof}

\begin{theorem}\label{thm:imaglb}
Let $n\in\N$, let $e_0,\dots,e_{\lfloor (n-1)/2\rfloor}\in\R$, put $e_k=0$ for
$k>\lfloor (n-1)/2\rfloor$, and $S=\sum_k e_k\,x^{n-1-2k}(y^2+t^2)^k$. Then the following are
equivalent:
\begin{enumerate}
  \item $v=y\,S$ satisfies $t\,\Delta v=\partial_t v$;
  \item $w=t\,S$ satisfies $t^2\Delta w-t\,\partial_t w+w=0$;
  \item $(2k+3)(2k+2)\,e_{k+1}=-(n-2k-1)(n-2k-2)\,e_k$ for all $k\ge0$.
\end{enumerate}
Consequently, for every $n\in\N$, the coefficient $v_n$ of $\ii$ in $z^n$ satisfies the same
Laplace--Beltrami equation as $\Real(z^n)$, the coefficient $w_n$ of $\jj$ satisfies
\[
  \boxed{\;t^2\Delta w_n-t\,\frac{\partial w_n}{\partial t}+w_n=0,\;}
\]
and $S_n$ is the unique polynomial of this form with $e_0=n$.
\end{theorem}

\begin{proof}
As in Theorem~\ref{thm:lb},
\[
  P_{xx}+4rP_{rr}+6P_r
  =\sum_{k\ge 0}\Bigl[(n-1-2k)(n-2-2k)\,e_k+(2k+2)(2k+3)\,e_{k+1}\Bigr]x^{n-3-2k}r^k,
\]
now using $4k(k-1)+6k=2k(2k+1)$. By Lemma~\ref{lem:axialim} each equation holds iff this
polynomial vanishes on $\{r>1\}$, that is iff every bracket vanishes. The remaining assertions
follow from Corollary~\ref{cor:srec}.
\end{proof}

\begin{remark}
The three equations say that $f=z^n=u+\ii v+\jj w$ has components $u,v$ in the kernel of the
hyperbolic Laplace--Beltrami operator $t^2\Delta-t\,\partial_t$ and $w$ in the kernel of
$t^2\Delta-t\,\partial_t+1$. Conceptually $\rho\,S(x,\rho^2)=\operatorname{Im}(x+\ii\rho)^n$
is harmonic in the $(x,\rho)$-plane, and $\rho\,(P_{xx}+4rP_{rr}+6P_r)=(\partial_x^2+\partial_\rho^2)(\rho P)$
for $r=\rho^2$.
\end{remark}

\section{Clifford algebras}

Nothing above used more than $v^2=-r$. Let $\mu\ge 1$ and let $\mathrm{Cl}_{0,\mu}$ be the real
Clifford algebra with generators $e_1,\dots,e_\mu$, $e_ie_j+e_je_i=-2\delta_{ij}$. A
\emph{paravector} is
\[
  z=x+v,\qquad v=\sum_{i=1}^{\mu}x_ie_i,\qquad r=\abs{v}^2=\sum_{i=1}^{\mu}x_i^2 ,
\]
and $\mu=2$, $(e_1,e_2)=(\ii,\jj)$, $(x_1,x_2)=(y,t)$ is the case treated so far. Write
$R_n(x,r)=\sum_kc_k\,x^{n-2k}r^k$ and $S_n(x,r)=\sum_ks_k\,x^{n-1-2k}r^k$ with the
coefficients $c_k=(-1)^k\binom{n}{2k}$, $s_k=(-1)^k\binom{n}{2k+1}$ of
Theorems~\ref{thm:closed} and~\ref{thm:imag}.

\begin{theorem}[Powers of paravectors]\label{thm:clifford}
Let $A$ be a real associative unital algebra, $x,r\in\R$ and $w\in A$ with $w^2=-r$. Then for
every $n\in\N$
\[
  \boxed{\;(x+w)^n=R_n(x,r)+S_n(x,r)\,w.\;}
\]
In particular, for every paravector $z=x+v$ of $\mathrm{Cl}_{0,\mu}$,
\[
  z^n=R_n\bigl(x,\abs{v}^2\bigr)+\sum_{i=1}^{\mu}x_i\,S_n\bigl(x,\abs{v}^2\bigr)\,e_i ,
\]
so powers of paravectors are paravectors, and the coefficients $c_k$, $s_k$, their recursions
and symmetries do not depend on $\mu$.
\end{theorem}

\begin{proof}
$x$ is central, so $(x+w)^n=\sum_i\binom{n}{i}x^{n-i}w^i$ with $w^{2k}=(-r)^k$ and
$w^{2k+1}=(-r)^kw$. For a paravector, $v^2=-\abs{v}^2$ by the defining relations.
\end{proof}

Let $\Delta=\partial_x^2+\sum_{i=1}^{\mu}\partial_{x_i}^2$ on $\R^{\mu+1}$ and fix an index $l$;
the coordinate $x_l$ plays the role of $t$.

\begin{lemma}\label{lem:axialN}
Let $P\in\R[x,r]$, $u=P(x,\abs{v}^2)$ and $g_i=x_i\,P(x,\abs{v}^2)$. Then, with all
polynomials evaluated at $r=\abs{v}^2$,
\begin{align*}
  x_l\,\Delta u-(\mu-1)\,\partial_{x_l}u&=x_l\,\bigl(P_{xx}+4rP_{rr}+2P_r\bigr),\\
  x_l\,\Delta g_i-(\mu-1)\,\partial_{x_l}g_i&=x_l\,x_i\,\bigl(P_{xx}+4rP_{rr}+6P_r\bigr)\qquad(i\ne l),\\
  x_l^2\,\Delta g_l-(\mu-1)\,x_l\,\partial_{x_l}g_l+(\mu-1)\,g_l&=x_l^3\,\bigl(P_{xx}+4rP_{rr}+6P_r\bigr).
\end{align*}
\end{lemma}

\begin{proof}
$\partial_{x_i}u=2x_iP_r$ and $\partial_{x_i}^2u=2P_r+4x_i^2P_{rr}$, hence
$\Delta u=P_{xx}+4rP_{rr}+2\mu P_r$; moreover $\Delta(x_iu)=x_i\,\Delta u+2\,\partial_{x_i}u$.
In each line the terms with $\mu$ cancel.
\end{proof}

\begin{theorem}\label{thm:cliffordlb}
Let $n\in\N$ and $\mu\ge1$.
\begin{enumerate}
  \item For $u=\sum_kd_k\,x^{n-2k}\abs{v}^{2k}$ the equation
  $x_l\,\Delta u=(\mu-1)\,\partial_{x_l}u$ holds iff
  $(2k+2)(2k+1)\,d_{k+1}=-(n-2k)(n-2k-1)\,d_k$ for all $k$.
  \item For $S=\sum_ke_k\,x^{n-1-2k}\abs{v}^{2k}$ and $i\ne l$ the equation
  $x_l\,\Delta(x_iS)=(\mu-1)\,\partial_{x_l}(x_iS)$ holds iff
  $(2k+3)(2k+2)\,e_{k+1}=-(n-2k-1)(n-2k-2)\,e_k$ for all $k$.
  \item For the same $S$ and $w=x_lS$ the equation
  $x_l^2\,\Delta w-(\mu-1)\,x_l\,\partial_{x_l}w+(\mu-1)\,w=0$ holds iff the recursion in (2) holds.
\end{enumerate}
Consequently, for $z^n=u_n+\sum_iv_{n,i}\,e_i$ in $\mathrm{Cl}_{0,\mu}$ and every $n$,
\[
  \boxed{\;
  \begin{aligned}
  x_l\,\Delta f&=(\mu-1)\,\frac{\partial f}{\partial x_l}
    &&\text{for } f\in\{u_n\}\cup\{v_{n,i}:i\ne l\},\\
  x_l^2\,\Delta w-(\mu-1)\,x_l\,\frac{\partial w}{\partial x_l}+(\mu-1)\,w&=0
    &&\text{for } w=v_{n,l}.
  \end{aligned}
  \;}
\]
\end{theorem}

\begin{proof}
By Lemma~\ref{lem:axialN} the three defects are the brackets of Theorems~\ref{thm:lb}
and~\ref{thm:imaglb}, evaluated at $r=\abs{v}^2$, times a monomial in $x_l,x_i$ that is non-zero
on an open set on which $r$ takes all values $>1$. (In (2) one needs $\mu\ge2$ for $i\ne l$ to exist.)
\end{proof}

\begin{remark}
For $\mu=2$, $l=2$ these are the equations of Sections~2 and~3. In general they say that the
components $u_n$, $v_{n,i}$ ($i\ne l$) lie in the kernel of the Laplace--Beltrami operator
$x_l^2\Delta-(\mu-1)x_l\partial_{x_l}$ of the hyperbolic metric $x_l^{-2}(dx^2+\sum dx_i^2)$ on
the upper half space of $\R^{\mu+1}$, and $v_{n,l}$ in the kernel of
$x_l^2\Delta-(\mu-1)x_l\partial_{x_l}+(\mu-1)$. The index $l$ is arbitrary: by symmetry of $z^n$
in the $x_i$ every coordinate may serve as the distinguished one.
\end{remark}

\section{Examples}

\begin{example}[The polynomials]\label{ex:table}
For $n\le 6$, with $r=y^2+t^2$ (or $r=\abs{v}^2$ in $\mathrm{Cl}_{0,\mu}$), $z^n=R_n+S_n\,v$ with
\begin{center}
\renewcommand{\arraystretch}{1.15}
\begin{tabular}{@{}cll@{}}
  $n$ & $R_n(x,r)=\Real(z^n)$ & $S_n(x,r)$\\\hline
  $1$ & $x$ & $1$\\
  $2$ & $x^2-r$ & $2x$\\
  $3$ & $x^3-3xr$ & $3x^2-r$\\
  $4$ & $x^4-6x^2r+r^2$ & $4x^3-4xr$\\
  $5$ & $x^5-10x^3r+5xr^2$ & $5x^4-10x^2r+r^2$\\
  $6$ & $x^6-15x^4r+15x^2r^2-r^3$ & $6x^5-20x^3r+6xr^2$
\end{tabular}
\end{center}
The even rows reproduce the examples of the notes; for $n=8$ and $n=10$ the coefficients
$a_k$ are $1,-28,70,-28,1$ and $1,-45,210,-210,45,-1$.
\end{example}

\begin{example}[Running the recursions, $n=6$]
Corollary~\ref{cor:rec} gives
\[
  c_0=1,\quad c_1=-\tfrac{6\cdot5}{2\cdot1}\,c_0=-15,\quad
  c_2=-\tfrac{4\cdot3}{4\cdot3}\,c_1=15,\quad
  c_3=-\tfrac{2\cdot1}{6\cdot5}\,c_2=-1,\quad c_4=-\tfrac{0\cdot(-1)}{8\cdot7}\,c_3=0,
\]
and Corollary~\ref{cor:srec} gives
\[
  s_0=6,\quad s_1=-\tfrac{5\cdot4}{3\cdot2}\,s_0=-20,\quad
  s_2=-\tfrac{3\cdot2}{5\cdot4}\,s_1=6,\quad s_3=-\tfrac{1\cdot0}{7\cdot6}\,s_2=0 .
\]
Both recursions terminate by themselves. The symmetries read $a_{3-k}=-a_k$ and
$s_{2-k}=s_k$ ($m=3$).
\end{example}

\begin{example}[A numerical power]
For $z=1+2\ii+3\jj$ one has $x=1$, $r=13$, hence $R_3=1-39=-38$, $S_3=3-13=-10$ and
\[
  z^3=-38-20\,\ii-30\,\jj .
\]
\end{example}

\begin{example}[The equations for $n=3$, $\mu=2$]
$u=x^3-3xr$ has $\Delta u=6x-12x=-6x$, so $t\,\Delta u=-6xt=\partial_tu$.
$v=y(3x^2-r)$ has $\Delta v=6y-\Delta(yr)=6y-8y=-2y$, so $t\,\Delta v=-2yt=\partial_tv$.
$w=t(3x^2-r)$ has $\Delta w=-2t$ and $\partial_tw=3x^2-r-2t^2$, so
$t^2\Delta w-t\,\partial_tw+w=-2t^3+2t^3=0$.
\end{example}

\begin{example}[The equations in $\mathrm{Cl}_{0,3}$, $n=2,3$, $l=3$]
Here $\mu-1=2$ and $r=x_1^2+x_2^2+x_3^2$. For $z^2=x^2-r+2x\,v$:
$u=x^2-r$ has $\Delta u=2-6=-4$, so $x_3\Delta u=-4x_3=2\,\partial_{x_3}u$; the components
$2xx_1$, $2xx_2$ are harmonic and independent of $x_3$; and $w=2xx_3$ is harmonic with
$-2x_3\,\partial_{x_3}w+2w=-4xx_3+4xx_3=0$.
For $z^3$: $u=x^3-3xr$ has $\Delta u=6x-18x=-12x$ and $2\,\partial_{x_3}u=-12xx_3=x_3\Delta u$.
The same $u$ in $\mathrm{Cl}_{0,2}$ has $\Delta u=-6x$: the polynomial is unchanged, the
Laplacian and the factor $\mu-1$ change together.
\end{example}

\section{Isolated roots outside the $xy$-plane}

Consider polynomials $p(q)=\sum_{n=0}^{N}q^n a_n$ with quaternionic coefficients on the right,
evaluated at reduced quaternions $z=x+\ii y+\jj t$. The $xy$-plane $\{t=0\}$ is
$\mathbb{C}_{\ii}=\{q:\ q_{\jj}=q_{\ii\jj}=0\}$. Is there a polynomial of degree $2$ with an
isolated root outside this plane? The answer depends on the coefficients.

\begin{theorem}[Coefficients in the $xy$-plane]\label{thm:circle}
Let all $a_n\in\mathbb{C}_{\ii}$ and let $z_0=x+\ii y+\jj t$ be a root of $p$ with $t\ne0$. Then
every $x+\ii y'+\jj t'$ with $y'^2+t'^2=y^2+t^2$ is a root of $p$. In particular $z_0$ is not
isolated: for every $\varepsilon>0$ there is a root $z_1\ne z_0$ with $\abs{z_1-z_0}<\varepsilon$.
This holds for every degree $N$.
\end{theorem}

\begin{proof}
By Theorem~\ref{thm:clifford}, $z^n=R_n+v\,S_n$ with $R_n,S_n$ real, hence
\[
  p(z)=A+v\,B,\qquad A=\sum_nR_n(x,r)\,a_n,\quad B=\sum_nS_n(x,r)\,a_n,
\]
where $A,B\in\mathbb{C}_{\ii}$ depend on $(x,r)$ only. For $B=B_0+B_1\ii$ one computes
$v\,B=-yB_1+yB_0\,\ii+tB_0\,\jj-tB_1\,\ii\jj$. The $\jj$- and $\ii\jj$-components of
$p(z_0)=0$ therefore give $tB_0=tB_1=0$, so $B=0$ and then $A=0$; both vanish on the whole
circle. For the last claim rotate $(y,t)$ by a small angle $\theta$: the squared
distance is $2r(1-\cos\theta)\le r\theta^2$ and it is positive for $0<\theta\le1$.
\end{proof}

\begin{corollary}[Real coefficients, degree $2$]\label{cor:realquad}
For $b,c\in\R$,
\[
  z^2+bz+c=0\iff x^2-r+bx+c=0\ \text{ and }\ (2x+b)\,y=(2x+b)\,t=0 .
\]
So the roots are the real roots of $x^2+bx+c$ and, if $4c>b^2$, the whole circle
$x=-b/2$, $y^2+t^2=c-b^2/4$; the isolated roots are real.
\end{corollary}

\begin{proof}
$z^2+bz+c=(x^2-r+bx+c)+(2x+b)\,v$ by Example~\ref{ex:table}.
\end{proof}

\begin{theorem}[A coefficient outside the plane]\label{thm:jpoly}
The polynomial
\[
  \boxed{\;p(q)=q^2-q\,(1+\jj)+\jj\;}
\]
has in all of $\HH$ exactly the two roots $q=1$ and $q=\jj$. Thus $\jj=(0,0,1)$ is an isolated
root of a polynomial of degree $2$ that lies outside the $xy$-plane.
\end{theorem}

\begin{proof}
For $q=a+b\,\ii+c\,\jj+d\,\ii\jj$ the four components of $p(q)=0$ are
\begin{align*}
  a^2-b^2-c^2-d^2-a+c&=0, & (2a-1)\,b+d&=0,\\
  2ac-c-a+1&=0, & (2a-1)\,d-b&=0 .
\end{align*}
The right column gives $b\,((2a-1)^2+1)=0$, so $b=d=0$. Then the first equation reads
$(a-c)(a+c-1)=0$. If $a=c$ the third becomes $2a^2-2a+1=0$, which has no real solution; if
$c=1-a$ it becomes $a(a-1)=0$, giving $q=\jj$ and $q=1$. Conversely both are roots, since
$\jj^2-\jj(1+\jj)+\jj=-1-\jj+1+\jj=0$.
\end{proof}

\begin{remark}
All coefficients of this $p$ lie in the plane $\mathbb{C}_{\jj}$, where $p(w)=(w-1)(w-\jj)$.
Theorem~\ref{thm:circle} holds for every plane $\mathbb{C}_I$ in place of $\mathbb{C}_{\ii}$: an
isolated root outside $\mathbb{C}_I$ needs a coefficient outside $\mathbb{C}_I$.
\end{remark}

\section{Degree three and general degree}

The answer of Section~6 does not depend on the degree. By Theorem~\ref{thm:circle} a polynomial
of any degree with coefficients in the $xy$-plane has no isolated root outside it. For real
coefficients the roots can be described completely.

\begin{theorem}[Real coefficients, any degree]\label{thm:realpoly}
Let $\alpha_0,\dots,\alpha_N\in\R$, $p(z)=\sum_nz^n\alpha_n$ and
\[
  A(x,r)=\sum_{n=0}^{N}R_n(x,r)\,\alpha_n,\qquad B(x,r)=\sum_{n=0}^{N}S_n(x,r)\,\alpha_n .
\]
Then $p(z)=A+B\,v$, and
\[
  p(z)=0\iff A(x,r)=0\ \text{ and }\ B(x,r)\,y=B(x,r)\,t=0 .
\]
So the roots are the real roots of $p$ and whole circles $\{x\}\times\{y^2+t^2=r\}$ with
$A(x,r)=B(x,r)=0$, $r>0$; the isolated roots are real.
\end{theorem}

\begin{proof}
Theorem~\ref{thm:clifford} and comparison of components; $v=\ii y+\jj t$.
\end{proof}

\begin{corollary}[Real coefficients, degree $3$]\label{cor:realcubic}
For $p(z)=z^3+bz^2+cz+d$ with $b,c,d\in\R$,
\[
  A=x^3-3xr+b\,(x^2-r)+cx+d,\qquad B=3x^2-r+2bx+c .
\]
A root with $(y,t)\ne(0,0)$ satisfies
\[
  \boxed{\;r=3x^2+2bx+c=p'(x),\qquad 8x^3+8bx^2+2(b^2+c)\,x+bc-d=0.\;}
\]
\end{corollary}

\begin{proof}
$B=0$ gives $r$; substituting into $A=0$ gives $-A=8x^3+8bx^2+2(b^2+c)x+bc-d$, precisely
$(3x+b)\,B-A$.
\end{proof}

To produce isolated roots outside the $xy$-plane one needs coefficients outside it. The general
tool is the analogue of Theorem~\ref{thm:circle} on all of $\HH$.

\begin{theorem}[Coefficients in $\mathbb{C}_{\jj}$, all of $\HH$]\label{thm:sphere}
Every quaternion is $q=x+w$ with $w=b\,\ii+c\,\jj+d\,\ii\jj$, $w^2=-\abs{w}^2$, hence
$q^n=R_n(x,\abs{w}^2)+S_n(x,\abs{w}^2)\,w$ and $p(q)=A+w\,B$ for every polynomial. If all
$a_n\in\mathbb{C}_{\jj}=\{a+c\,\jj\}$ and $q$ is a root with $(b,d)\ne(0,0)$, then $A=B=0$ and
every $x+w'$ with $\abs{w'}=\abs{w}$ is a root.
\end{theorem}

\begin{proof}
For $B=B_0+B_2\,\jj$ the $\ii$- and $\ii\jj$-components of $w\,B$ are $bB_0-dB_2$ and
$bB_2+dB_0$, while $A$ has none. So $(b^2+d^2)B_0=(b^2+d^2)B_2=0$.
\end{proof}

\begin{theorem}[Degree $3$]\label{thm:jcubic}
The polynomial
\[
  \boxed{\;p(q)=q^3-q^2(1+3\jj)+q\,(-2+3\jj)+2\;}
\]
has in all of $\HH$ exactly the three roots $1$, $\jj$ and $2\jj$. All are isolated, and
$\jj=(0,0,1)$ and $2\jj=(0,0,2)$ lie outside the $xy$-plane.
\end{theorem}

\begin{proof}
On $\mathbb{C}_{\jj}$ everything commutes and $p(q)=(q-1)(q-\jj)(q-2\jj)$; $\HH$ has no zero
divisors, so the roots in $\mathbb{C}_{\jj}$ are $1,\jj,2\jj$. A root outside $\mathbb{C}_{\jj}$
would by Theorem~\ref{thm:sphere} force both $x+\jj\rho$ and $x-\jj\rho$, $\rho>0$, to be
roots, but $-\rho\notin\{0,1,2\}$.
\end{proof}

\begin{theorem}[Degree $n$]\label{thm:jpow}
Let $n\ge1$ and $p(q)=q^n-\jj$.
\begin{enumerate}
  \item $p$ has a root, for instance $q=\cos\frac{\pi}{2n}+\jj\,\sin\frac{\pi}{2n}$.
  \item Every root lies in $\mathbb{C}_{\jj}$ and has $t\ne0$; so \emph{all} roots lie outside the
  $xy$-plane.
  \item $p$ has only finitely many roots in $\HH$; so every root is isolated.
\end{enumerate}
Hence for every degree $n\ge1$ there is a polynomial of degree $n$ all of whose roots are
isolated and lie outside the $xy$-plane.
\end{theorem}

\begin{proof}
(2) $q^n=R_n+S_n\,w=\jj$ gives $S_nb=S_nd=0$ and $S_nc=1$, so $S_n\ne0$, $b=d=0$ and
$c\ne0$. (1), (3) The map $a+c\,\ii\mapsto a+c\,\jj$ is an injective ring homomorphism
$\mathbb{C}\to\HH$ onto $\mathbb{C}_{\jj}$; by (2) it maps the finitely many solutions of
$w^n=\ii$ in $\mathbb{C}$ onto the roots of $p$, and $w=e^{\ii\pi/(2n)}$ is one of them.
\end{proof}

\begin{remark}
A real cubic has three real roots, or one real root and one circle: the circle over the real
root $x$ of the resolvent in Corollary~\ref{cor:realcubic} with $p'(x)>0$ is the pair of
complex conjugate roots $x\pm\ii\sqrt{p'(x)}$ rotated about the $x$-axis. In Theorem~\ref{thm:jpow}
the number of roots is exactly $n$. These two counts are classical and not formalised here.
\end{remark}

\section{Formal verification}

The file \texttt{RealPartsOfPowers.lean} (Lean~4, toolchain \texttt{v4.34.0}, Mathlib
\texttt{v4.34.0}) works in Mathlib's quaternions $\HH[\R]$ with $z=\langle x,y,t,0\rangle$ and
verifies every statement of this paper:

{\small
\begin{longtable}{@{}>{\raggedright\arraybackslash}p{0.42\textwidth}>{\raggedright\arraybackslash}p{0.54\textwidth}@{}}
  \lean{re\_pow}, \lean{re\_pow\_even} & Theorem~\ref{thm:closed}, every $n$, and its form for $n=2m$\\
  \lean{c\_zero}, \lean{c\_succ}, \lean{a\_zero}, \lean{a\_succ} & Corollary~\ref{cor:rec} (over $\mathbb{Z}$, division-free, all $k$)\\
  \lean{a\_self}, \lean{a\_sub} & Corollary~\ref{cor:sym}\\
  \lean{laplaceBeltrami\_defect\_general} & Lemma~\ref{lem:axial}, every polynomial $P$\\
  \lean{laplaceBeltrami\_defect} & the coefficient formula in the proof of Theorem~\ref{thm:lb}\\
  \lean{laplaceBeltrami\_iff} & Theorem~\ref{thm:lb}, the equivalence\\
  \lean{laplaceBeltrami\_re\_pow} & Theorem~\ref{thm:lb}, $t\,\Delta u_n=\partial_tu_n$ for every $n$\\
  \lean{eq\_c\_of\_recursion} & Theorem~\ref{thm:lb}, uniqueness\\
  \lean{imI\_pow}, \lean{imJ\_pow}, \lean{imK\_pow} & Theorem~\ref{thm:imag}\\
  \lean{s\_zero}, \lean{s\_succ}, \lean{s\_sub} & Corollary~\ref{cor:srec}\\
  \lean{defect\_VI}, \lean{defect\_VJ} & Lemma~\ref{lem:axialim}, every polynomial $P$\\
  \lean{laplaceBeltrami\_imI\_iff} & Theorem~\ref{thm:imaglb}, (1)\,$\Leftrightarrow$\,(3)\\
  \lean{hyperbolic\_imJ\_iff} & Theorem~\ref{thm:imaglb}, (2)\,$\Leftrightarrow$\,(3)\\
  \lean{laplaceBeltrami\_imI\_pow} & Theorem~\ref{thm:imaglb}, the equation for $v_n$\\
  \lean{hyperbolic\_imJ\_pow} & Theorem~\ref{thm:imaglb}, the equation for $w_n$\\
  \lean{eq\_s\_of\_recursion} & Theorem~\ref{thm:imaglb}, uniqueness\\
  \lean{pow\_eq\_of\_sq}, \lean{z\_pow\_eq} & Theorem~\ref{thm:clifford}, any algebra; the quaternions as instance\\
  \lean{clifford\_pow}, \lean{clifford\_pow\_components} & Theorem~\ref{thm:clifford} in $\mathrm{Cl}_{0,\mu}$\\
  \lean{defectN\_re}, \lean{defectN\_im}, \lean{defectN\_last} & Lemma~\ref{lem:axialN}, every polynomial $P$\\
  \lean{hyperbolic\_re\_iff}, \lean{hyperbolic\_im\_iff}, \lean{hyperbolic\_last\_iff} & Theorem~\ref{thm:cliffordlb}, (1)--(3)\\
  \lean{hyperbolic\_scalar\_pow}, \lean{hyperbolic\_vector\_pow}, \lean{hyperbolic\_last\_pow} & Theorem~\ref{thm:cliffordlb}, the equations for $z^n$\\
  \lean{Rpoly\_examples}, \lean{Spoly\_examples}, \lean{a\_four}, \lean{a\_five}, \lean{s\_six} & Example~\ref{ex:table} and the coefficient lists\\
  \lean{z\_cube\_example} & $(1+2\ii+3\jj)^3=-38-20\ii-30\jj$\\
  \lean{evalPoly\_z}, \lean{polyA\_polyB\_eq\_zero} & $p(z)=A+vB$ and $A=B=0$ in Theorem~\ref{thm:circle}\\
  \lean{circle\_of\_roots}, \lean{root\_not\_isolated} & Theorem~\ref{thm:circle}\\
  \lean{real\_quadratic\_root\_iff} & Corollary~\ref{cor:realquad}, the equivalence\\
  \lean{jPoly\_root\_iff}, \lean{jPoly\_root\_outside\_plane} & Theorem~\ref{thm:jpoly}\\
  \lean{real\_poly\_root\_iff} & Theorem~\ref{thm:realpoly}, the equivalence\\
  \lean{real\_cubic\_root\_iff}, \lean{real\_cubic\_circle} & Corollary~\ref{cor:realcubic}\\
  \lean{quat\_pow\_eq}, \lean{evalPoly\_quat}, \lean{polyA\_polyB\_eq\_zero\_J}, \lean{sphere\_of\_roots\_J} & Theorem~\ref{thm:sphere}\\
  \lean{jCubic\_root\_iff} & Theorem~\ref{thm:jcubic}\\
  \lean{exists\_pow\_eq\_j}, \lean{pow\_eq\_j}, \lean{pow\_eq\_j\_finite} & Theorem~\ref{thm:jpow}, (1)--(3)
\end{longtable}}

The formal proof of Theorem~\ref{thm:closed} follows the proof above verbatim:
\lean{Commute.add\_pow}, a parity split of $\sum_{i\le n}$, and $v^{2k}=(-r)^k$.
For Lemma~\ref{lem:axial} a polynomial is a finite sum of monomials $d_i\,x^{p_i}r^{k_i}$,
partial derivatives are Mathlib's \lean{deriv} in one coordinate, and
$\Delta$ is the sum of the three second partial derivatives. The direction
``equation $\Rightarrow$ recursion'' of Theorem~\ref{thm:lb} evaluates at $x=1$, $y=0$,
$t=\sqrt r$ and uses that a real polynomial with infinitely many roots vanishes.
The Clifford algebra is Mathlib's \lean{CliffordAlgebra} of the quadratic form
$-\sum x_i^2$ on $\R^\mu$, $e_i$ is the image of the $i$-th unit vector, and
partial derivatives in $x_i$ are \lean{deriv} along \lean{Function.update}. The component
formula is proved as an identity in the algebra; linear independence of $1,e_1,\dots,e_\mu$,
which makes the components unique, is not used and not formalised.
All results depend only on the axioms \lean{propext}, \lean{Classical.choice},
\lean{Quot.sound}; the build gate fails on any other axiom, in particular on \lean{sorryAx}.

\end{document}
